\documentclass[10pt,a4paper]{amsart}
\usepackage[margin=19mm]{geometry}
\usepackage{amsmath,amssymb,mathtools}
\usepackage[scr=rsfs,cal=euler]{mathalfa}
\usepackage{xfrac}
\usepackage[unicode,hidelinks,bookmarks=false]{hyperref}
\newcommand{\ie}{i.e.\ }
\newcommand{\cf}{cf.\ }
\newcommand{\resp}{respectively\ }
\newcommand{\Subsection}{\subsection}
\providecommand{\nobreakdash}{\nobreak\hbox{-}}
\setlength{\emergencystretch}{2em}
\multlinegap0pt
\usepackage{bm}
\usepackage{bbm}
\usepackage{dsfont}
\usepackage{xspace}
\usepackage{cancel}
\usepackage{mathtools}
\usepackage{amsthm}
\makeatletter
\@ifundefined{thm}{\newtheorem{thm}{Theorem}}{}
\makeatother
\def\BB{{\mathbb B}}
\def\CD{{\mathcal D}}
\def\CV{{\mathcal V}}
\def\fD{{\mathfrak D}}
\def\fa{{\mathfrak a}}
\def\fb{{\mathfrak b}}
\def\fc{{\mathfrak c}}
\def\ft{{\mathfrak t}}
\def\f{\frac}
\def\k{{\kappa}}
\def\sE{{\mathsf E}}
\def\sG{{\mathsf G}}
\def\sQ{{\mathsf Q}}
\def\sW{{\mathsf W}}
\title[Approximation and orthogonality on fully symmetric domains]{Approximation and orthogonality on fully symmetric domains}
\author{Yuan Xu}
\date{Source: arXiv:2509.10342v1 (CC BY 4.0)}
\begin{document}
\maketitle

\noindent\emph{Source and licence.} Approximation and orthogonality on fully symmetric domains. Version arXiv:2509.10342v1 (CC BY 4.0), distributed under the Creative Commons Attribution 4.0 International licence, as recorded in the arXiv licence metadata for this exact version and shown on the abstract page https://arxiv.org/abs/2509.10342v1. Full rights notice: licenses/arxiv-2509.10342-CC-BY-4.0.txt.\par\smallskip

\noindent\textit{Editorial scope.} Counted unit: the Thm whose source label is \texttt{thm:PDE\_d=2} in the article (source0.tex lines 862--880, source label \texttt{thm:PDE\_d=2}), delivered together with the article's own complete original proof of it (source0.tex lines 882--915, one \texttt{proof} environment of 34 lines). No mathematics is written, repaired or re-derived: statement and proof are the article's own text line for line. Required context retained uncounted: eq:eigenB2 (lines 423--425); sQ-sG (lines 845--849). Every definition, display and label of those lines is retained unchanged. Omitted from this package and not redistributed: 1--422, 426--844, 850--861, 881--881, 916--1819. Nothing inside a retained excerpt is omitted or altered: every retained body is delivered exactly as the article's lines are written, so no omission receipt of any kind accompanies this package. Citations: the retained excerpts carry no citation command at all, so no bibliography entry of the article is transcribed and no reference is redistributed. Reference adaptation: none was needed and none was made; every reference inside the delivered unit resolves inside it, so no unresolved reference or citation is produced. Printed slips kept literal: no printed word, symbol, hypothesis or number of the counted statement or of its proof was altered. Layout and class: the article's own class is \texttt{amsart} with packages including amsmath, amsthm, amsfonts, amssymb, accents, enumerate, color, graphicx, pict2e; the wrapper is the lane's pinned \texttt{amsart} editorial wrapper at 10pt on A4, which restates the theorem-like environments the retained text needs and adds the article's own macro definitions that the retained text needs (\texttt{\textbackslash BB}, \texttt{\textbackslash CD}, \texttt{\textbackslash CV}, \texttt{\textbackslash f}, \texttt{\textbackslash fD}, \texttt{\textbackslash fa}, \texttt{\textbackslash fb}, \texttt{\textbackslash fc}, \texttt{\textbackslash ft}, \texttt{\textbackslash k}, \texttt{\textbackslash sE}, \texttt{\textbackslash sG}, \texttt{\textbackslash sQ}, \texttt{\textbackslash sW}), with extra wrapper packages (bm, bbm, dsfont, xspace, cancel, mathtools). The delivered numbering is the unit's own. Assets: no retained span contains a figure environment, a graphics-inclusion command, a PDF-page inclusion, a table or a TikZ picture; no image file is copied into this package, the package declares no asset dependency and no image file is redistributed with it. Licence: version arXiv:2509.10342v1 of this article is distributed under the Creative Commons Attribution 4.0 International licence, as recorded in the arXiv licence metadata for this exact version and shown on the abstract page https://arxiv.org/abs/2509.10342v1. A full rights notice is delivered at licenses/arxiv-2509.10342-CC-BY-4.0.txt. Characters: every retained excerpt is pure ASCII, so the delivered engine, XeLaTeX with its default Latin Modern text font, reports no missing character.
\begin{equation}\label{eq:eigenB2}
   \CD^\k_\BB Y = - n (n+ 2 |\k| + 2) Y, \qquad Y\in \CV_n(\sW^\k_\BB, \BB^2).
\end{equation}

\begin{align} \label{sQ-sG}
   \sQ_{j,n}^\k(u,v)\, &  = C_{n-2j}^{(2j + \k_2+\k_3 +1,\k_1)}(u) (1-u^2)^{j} 
   P_j^{(\k_3, \k_2-\f12)}\bigg(2 \frac{-\fa+(\fa-\fc)u^2+v^2}{(\fb-\fa)(1-u^2)} -1\bigg) \notag  \\
   \,& = (\sG_{2j,n}^\k \circ \psi) (u,v).
\end{align}

\begin{thm}\label{thm:PDE_d=2}
For $\k_1 =0$, the polynomials in $\CV_n^{\circ, \sE}(\sW_{\fa,\fb,\fc}^\k, \Omega_{\fa,\fb,\fc})$ 
are eigenfunctions of a differential operator, 
\begin{equation} \label{eq:spectral}
    \fD_{\fa,\fb,\fc}^\k Y = -n(n+ 2\k_2+2\k_3+ 2) Y, \qquad Y \in \CV_n^{\circ, \sE}(\sW_{\fa,\fb,\fc}^\k, \Omega_{\fa,\fb,\fc}),
\end{equation} 
and the operator is given explicitly by, for $(u,v) \in \Omega_{\fa,\fb,\fc}$, 
\begin{align*}
 \fD_{\fa,\fb,\fc}^\k \, & = (1-u^2) \partial_{uu} - 2(-\fc+v^2) \frac{u}{v} \partial_{u v} + \left[ \fa + \fb - v^2
       + \frac{- \fa \fb + (\fa - \fc) (\fb - \fc) u^2}{v^2}\right]\partial_{vv}\\
  & + \left[ -2\fa + \fc -\frac{- \fa \fb + (\fa - \fc) (\fb - \fc) u^2}{v^2} + \right]\frac{1}{v}\partial_{v} 
    + (2|\k|+3) \left[ u \partial_u + (v^2-\fa)\frac{1}{v}\partial_{v}\right]\\
  &  + 2 \k_2 (\fb-\fa) \frac{1}{v} \partial_v. 
\end{align*}
Moreover, the operator is related to the second-order differential operator $\CD_{\BB}^\k$ by 
\begin{equation} \label{eq:fD=D}
        \fD_{\fa,\fb,\fc}^\k (f\circ \psi) = \left(\CD_\BB^\k f \right) \circ \psi.   
\end{equation}
\end{thm} 

\begin{proof}
Our goal is to find an operator $\fD_{\fa,\fb,\fc}^\k$ that satisfies \eqref{eq:fD=D}. 
Indeed, by \eqref{sQ-sG}, such an operator would imply 
\begin{align*}
  \fD_{\fa,\fb,\fc}^\k \sQ_{j,n}^\k \, & =  \fD_{\fa,\fb,\fc}^\k (\sG_{2j,n}^\k \circ \psi) 
       =  \left(\CD_\BB^\k \sG_{2j,n}^\k \right) \circ \psi \\ 
      &  = -n(n+2|\k|+2) \sG_{2j,n}^\k \circ \psi =  -n(n+2|\k|+2) \sQ_{j,n}^\k,  
\end{align*}
where the second identity uses \eqref{eq:fD=D} and the third one follows from \eqref{eq:eigenB2},
which would have verified \eqref{eq:spectral}. 

Let $F(u,v) = (f\circ \psi)(u,v) = f(u, \ft(u,v))$. Denote by $\partial_1f$ and $\partial_2f$ the derivative of 
$f$ for the first and the second variable. Taking the derivative of $\partial_u F$ and $\partial_v F$
by chain rule and solve for $\partial_1 f$ and $\partial_2 f$ accordingly, it is straightforward to verify that
\begin{align*}
  \partial_2 f(u,\ft(u,v)) \, &= (\fb - \fa) \ft(u,v) \frac{1}{v} \partial_v F(u,v), \\
  \partial_1 f(u, \ft(u,v)) \, &= \partial_u F(u,v) - (\fa - \fc) \frac{u}{v} \partial_v F(u,v).
\end{align*}
Moreover, taking the derivatives one more time and solving for the second-order derivatives of $f$, we
obtain 
\begin{align*}
  \partial_2^2 f(u, t(u,v)) = \, &  (\fb-\fa)^2 \f{\ft(u,v)^2}{v^2} \partial_u^2 F(u,v) - (\fb - \fa)\big(-\fa+ (\fa - \fc)v^2\big)
   \frac{u}{v} \partial_v F(u,v), \\
  \partial_1 \partial_2 f(u, \ft(u,v))  = \, & (\fb - \fa) \frac{\ft(u,v)}{v} \partial_u \partial_v F(u,v)
         + (\fa - \fc)^2  \frac{u^2}{v^2} \partial^2_v F(u,v)\\
             \, &  - (\fa - \fc)(\fb-\fa)  \frac{u \ft(u,v)}{v^3} \partial_v F(u,v), \\
    \partial_1^2 f(u, \ft(u,v)) = \, & \partial_u^2 F(u,v) -2 (\fa - \fc) \frac{u}{v} \partial_u \partial_v F(u,v) 
        + (\fa-\fc)^2  \frac{u^2}{v^2} \partial^2_v F(u,v) \\
     & +  (\fa - \fc) \big((\fa - \fc) u^2 + v^2\big) \frac{1}{v^3} \partial_v F(u,v).       
\end{align*}
Putting these together and substituting them into $\CD_\BB^\k$ that has the coefficients, in front of the differentials,
in the variable $(s,t) = (u,t(u,v)))$, we obtain the stated expression for $\fD_{\fa,\fb,\fc}^\k$ after simplification by
\eqref{eq:fD=D}.  
\end{proof}

\end{document}
